\loadgeometry{narrow}%
\chapter*{第一瓶中之信}

{\setlength{\parindent}{0pt}\kaiti%
你好，世界。\\
再见，世界。\par
\halfline%
我能吞下痛苦而不伤身体，\\
却不能使一颗心免于悲伤。\par
\halfline%
即使引导已经破碎，\\
也请你继续前行。\par
\hfill%
雪莉}
\restoregeometry

\chapter{瓶装快乐}
“吵死了！”狂吠着、呐喊着、尖叫着、号哭着的手机被狠狠地掷出被窝；然而却不闻其落地的声音。“它不会被我掷出窗外了吧？那样的话，希望楼下有市议员能够接住它……”铃没有被不同于平日的、尖锐刺耳的闹钟叫醒，却被这样的骇人的想法惊醒。她本能地坐起，如同听闻北伐成功的老翁。手机静静地躺在咫尺之外的地板上，失去了发出所有声音的生气。铃很想确认它是死是活，但从被窝到床外地面的距离对她来说如同从织女星到牛郎星（相对性原理！）。她放弃了这样徒劳的冲动，重新把自己塞进被窝，用被子蒙住全身，仿佛这个无光的被窝就是她的全部表象世界。“这样看来，说不定那位科尼斯堡的幽灵也是个‘回笼教信徒’。”手机现在是否损坏对铃来说已经是超验的了。铃只能做好最坏的打算：“说不定也是好事；终于可以换新手机了；我一定要买最新款。”

\begin{branch}
\goodend 不知过了多久，铃终于恋恋不舍地告别了作为世界最后的温暖的被窝。她的脚刚踏上地面，就重重地踩在了本就已经残破的受害者上，使它一命呜呼。铃痛心地拾起它的尸体，回忆起它这么多年任劳任怨地被自己使用，流下了哀伤的泪水。“我一定要为它举办一场体面的葬礼。”葬礼结束后，铃将这位伟大的手机埋进它的预定的坟冢。她在墓前默哀，突然想到这位逝者的电子元器件构成的身体（含有铅、镉、汞、锂、塑料、溴化物等）会对土壤环境造成不可逆的破坏。铃于是不得不强忍着泪水将它的棺椁取出，重新摆放到出租屋中为其树立的灵牌前。在这场肃穆的葬礼期间，没有手机的铃丢失了所有的通讯消息；不过这无关紧要。告别旧友之后，铃重新购买了最新款的手机（分期！），回到了正常的生活之中。可喜可贺！
\end{branch}

可惜的是手机并没有去世，因为不久之后，它传来激昂的铃声（是手机传来铃声，不是铃传来铃声），像是整个红旗歌舞团在铃的卧室里演唱《货郎》。铃祈祷着手机的另一边赶快察觉到无人接听，因为她的被窝仅仅能遮蔽透过厚厚云层和卧室窗户的阳光，而不能阻断无形而刺耳的声音；铃的现象界出现了裂隙。铃尝试说服自己，听到铃声不过是她谬误的信念；这样相信着，铃声果然停止了。“希望那位数学家不要再设定什么欺人的恶魔了。”铃身体放松下来，像一块蜡似的融化在被窝里，“看来我的本质和蜡并无区别。”铃就快要复归梦乡。

然而，铃声又响起了，这次仿佛更加地热烈。“Scheiße!”铃徒劳地捂住自己的耳朵，铃声却不经中介似的进入铃的意识。“不会是大老板发现我没去研讨会了吧？”铃想象着最坏的可能。

\begin{branch}
\badend 铃立刻被这样的最大最恶的设想吓到了。她掀翻被子，跃下床，

\goodend 铃的腿立刻折了；这其一是因为刚睡醒的铃忘记了她的床是上床下桌的结构，她这是从一米多高的地方跳下来了；其二是因为铃从不运动导致身体很脆弱。于是铃在医院里度过了一段无比快乐的时光。同事们纷纷前来探望，假意地嘘寒问暖；不过这不重要。一想到自己还能合法合理地逃离工作那么长的时间，铃就发自内心地高兴。

\badend 铃慌忙抓起手机。来电人正是大老板。铃战战兢兢地按下接听键。

\goodend “铃！中了！我们中了！”电话那头的大老板压制着兴奋，“而且还是Oral！”铃感到一阵幸福的眩晕，“噫，我中了！”

\badend “有人和我说你没去我昨天叫你去的研讨会。”电话那头的大老板压制着愤怒，“你也没来实验室。你去哪了？是不想来了吗？”铃汗流浃背。“作为惩罚，这个月整理实验室和添置试剂都由你负责。”这下惨了。
\end{branch}

“不太可能。按照时间算，现在会议开始也没多久。即使被发现了，也不会这么快。”铃安抚着自己，更加紧地堵住耳朵。铃声终于再次停止。“肯定是学姐打电话来问我要数据。诶，她周末还干什么活。要是我能变成卡夫卡的甲虫就好了；毕竟甲虫不用工作。”

\begin{branch}
\goodend 铃这样想着，竟真感到自己的胸腹部长出了六条腿。她想要坐起查看自己的情况，却发现自己仅仅是在微微蜷曲着身体，而无法使上劲；这下铃真的变成甲虫了。“好耶！”铃高兴地用六只手拍了拍自己的肚子，发出清脆的响声。这下铃是彻底不用去上班了。至于人变作的甲虫应不应该参与社会保障体系，这个问题就留给法学家、政策学者和社会学家吧。
\end{branch}

当铃声再次响起，铃的半梦半醒的意识被击穿了。所有的可能性已经化作单纯的规定性，因为能够连着打来三次电话的，铃所知的仅仅有唯一的一位神人。铃不情愿地磨蹭着爬起、下床，接起了电话；因为她知道如果不接，对方还会一直无间断轰炸。电话的另一边传来熟悉的声音。

“铃！今天我休息。我们去逛街吧！”且不论人与人的休息日并不相通，铃今天仅仅是翘掉会议在赖床罢了。

“你找艾莉丝。”铃提出替代方案，祈求对方千万不要指名自己。

“艾莉丝说她在参加研讨会。”铃走出她的卧室环顾了一圈，才发现她的室友艾莉丝早就出门了。这也难怪，毕竟现在窗外的太阳（如果有）已经非常高了。“我可不用上早朝。”

“那你等她回来吧，”铃回忆着会议的日程，“她不久就会回来的，除非她掉到兔子洞里。”

“可是那就已经到下午了；我今天还要逛好多商场，买好多东西！”对方接着兴致勃勃地介绍着，仿佛在讲她的登陆作战计划。

“好吧好吧。”铃妥协了，赶紧打断这看起来无止境的话语，“在哪里会合？”当然，这里的“会合”指的不是作战旅。

“在你的门口哦！”对方说出意料之外的地点。接着门口传来着急的敲门声。自己什么时候把住址告诉她的？也许是艾莉丝告诉她的吧。

现在计算敲门声的拍子数是不合时宜的。铃慢慢走过去，取下防盗链，打开了门。

\begin{branch}
\badend 然而站在门外的人并不是铃设想的。也对，电话中的话语和敲门声即使有关联，也不一定是因果关系。并且，谁能界定那句标识位置的话是真实的还是仅仅是玩笑？门外的歹徒看到铃不按剧本地在家，慌了神。他在刹那间权衡了逃跑和行凶的利弊。也许是肾上腺素冲昏了他的头脑，他毫无理智地选择了后者。随着锐利的匕首刺进铃的腹部，铃倒在了血泊里。“该邀请谁主持我的葬礼呢？”铃开始构想葬礼的布置……
\end{branch}

门外站着的正是铃设想的那位人物。铃没好气地向她打招呼：“真早啊，彩名。”言外之意是“你影响到我睡懒觉了。”

“魔法少女莉莉姆，登场！”
